\chapter{Modular Forms and $L$-functions}
\begin{center}\itshape
Count the points of $y^2+y=x^3-x^2$ modulo $p$.\\
Expand $q\prod_{n\ge1}(1-q^n)^2(1-q^{11n})^2$.\\
The same numbers come out. Why?
\end{center}

In September 1955, at a symposium on algebraic number theory held in Tokyo and Nikk\=o, the young Japanese mathematician Yutaka Taniyama handed out a
list of problems. Problem 12 asked, in a slightly different form from the one known today, whether the $L$-function of every elliptic curve over $\Q$ comes from a
modular form. The question looked like wishful thinking. An elliptic curve is algebra: a cubic equation and its points. A modular form is analysis: a holomorphic
function on the upper half-plane with an infinite group of symmetries. Yet there was already evidence: Martin Eichler had shown in 1954 that, for modular forms like the product in the epigraph (two copies of Dedekind's
$\eta$ function), the coefficients count points modulo $p$ on a cubic curve.

Goro Shimura sharpened the question in the years after, and Andr\'e Weil showed in 1967 that it would follow from functional equations of $L$-functions, which made
it precise enough to test. Then, in 1985, Gerhard Frey noticed that a solution of $a^p+b^p=c^p$ would give an elliptic curve so strange that it could not be modular.
Ken Ribet proved in 1986 that it indeed could not be, and Andrew Wiles, with Richard Taylor, proved in 1995 that enough elliptic curves are modular to rule it out:
Fermat's Last Theorem fell to the conjecture of Taniyama's problem 12. The full statement, for every elliptic curve over $\Q$, was completed by Breuil, Conrad,
Diamond and Taylor in 2001.

Volume~5 met modular forms of level $1$: Eisenstein series, $\Delta$, its product formula, dimensions and the Hecke operators $T_p$. This chapter finishes the
proofs promised there (the fundamental domain), proves the valence formula and its consequence that $j$ classifies elliptic curves over $\C$, builds the whole
family of Hecke operators $T_n$, and proves Hecke's theorem: the $L$-function of a cusp form continues to the whole plane with a functional equation. The same
proof, applied to the product in the epigraph, gives the functional equation of the $L$-function of an elliptic curve, once that curve is known to be modular.

\begin{reminder}[from Volume 5, Chapter 12]
$\mathrm{SL}_2(\Z)$ acts on $\mathbb H$ and is generated by $T=\bigl(\begin{smallmatrix}1&1\\0&1\end{smallmatrix}\bigr)$ and $S=\bigl(\begin{smallmatrix}0&-1\\1&0\end{smallmatrix}\bigr)$
(Volume~5, Proposition~12.1); every orbit meets $F=\{|\operatorname{Re}\tau|\le\frac12,\ |\tau|\ge1\}$ (Theorem~12.2). $M_k$ and $S_k$ are the modular forms and cusp
forms of weight $k$, with $q=e^{2\pi i\tau}$; for $\alpha$ with $\det\alpha>0$, $(f|_k\alpha)(\tau)=(\det\alpha)^{k-1}(c\tau+d)^{-k}f(\alpha\tau)$. The Eisenstein series
$E_k=1-\frac{2k}{B_k}\sum\sigma_{k-1}(n)q^n$ lie in $M_k$ (Theorem~12.6); $\eta(\tau)=e^{\pi i\tau/12}\prod(1-q^n)$ satisfies $\eta(-1/\tau)=\sqrt{-i\tau}\,\eta(\tau)$, and
$\Delta=\eta^{24}\in S_{12}$ has no zeros on $\mathbb H$ (Theorem~12.7). $M_k=\C E_k\oplus\Delta M_{k-12}$ for $k\ge4$, $M_0=\C$, $M_2=0$ (Theorem~12.9), and
$1728\Delta=E_4^3-E_6^2$ (Corollary~12.11). $\rho=e^{2\pi i/3}$. \emph{Volume~4.} $\Gamma(s)=\int_0^\infty e^{-t}t^{s-1}\,\mathrm dt$ extends meromorphically to $\C$ with poles at
$0,-1,-2,\dots$ and no zeros (Volume~4, Theorem~9.36); Euler's pentagonal theorem $\prod_{n\ge1}(1-q^n)=\sum_{j\in\Z}(-1)^jq^{j(3j-1)/2}$ (Volume~4, Theorem~2.12).
\end{reminder}

\section{The fundamental domain}

\begin{theorem}\label{thm:funddomain}
\begin{enumerate}[label=(\alph*)]
\item If two distinct points $\tau,\tau'$ of $F$ lie in the same orbit of $\mathrm{SL}_2(\Z)$, then either $\operatorname{Re}\tau=\pm\frac12$ and $\tau'=\tau\mp1$, or
$|\tau|=1$ and $\tau'=-1/\tau$.
\item The stabiliser of $\tau\in F$ in $\mathrm{SL}_2(\Z)$ is $\{\pm I\}$, except for three points: $i$, with stabiliser $\{\pm I,\pm S\}$ of order $4$; $\rho$, with
the stabiliser of order $6$ generated by $ST$; and $-\overline\rho=\rho+1$, with the stabiliser of order $6$ generated by $TS$.
\end{enumerate}
\end{theorem}
\begin{proof}
Let $\tau\in F$ and $\gamma=\bigl(\begin{smallmatrix}a&b\\c&d\end{smallmatrix}\bigr)$ with $\gamma\tau\in F$. Exchanging $(\tau,\gamma)$ with $(\gamma\tau,\gamma^{-1})$ if
necessary, we may assume $\operatorname{Im}\gamma\tau\ge\operatorname{Im}\tau$, that is, $|c\tau+d|\le1$. Then $|c|\operatorname{Im}\tau\le1$ and
$\operatorname{Im}\tau\ge\frac{\sqrt3}2$, so $c\in\{0,\pm1\}$; replacing $\gamma$ by $-\gamma$, which acts in the same way, we may take $c\in\{0,1\}$.

If $c=0$, then $ad=1$ and $\gamma\tau=\tau\pm b$. Both real parts lie in $[-\frac12,\frac12]$, so $b=0$ and $\gamma=\pm I$, or $|b|=1$ and
$\operatorname{Re}\tau=\mp\frac12$.

If $c=1$, then $|\tau+d|\le1$. Now $|\tau+d|^2=|\tau|^2+2d\operatorname{Re}\tau+d^2\ge1-|d|+d^2$, which exceeds $1$ when $|d|\ge2$. If $d=1$, equality forces
$|\tau|=1$ and $\operatorname{Re}\tau=-\frac12$: $\tau=\rho$; if $d=-1$, similarly $\tau=-\overline\rho$. If $d=0$, then $|\tau|\le1$, so $|\tau|=1$, and $b=-1$ (the
determinant is $-b$), so $\gamma\tau=a-\frac1\tau=a-\overline\tau$. Since $-\overline\tau\in F$, this lies in $F$ only for $a=0$, or for $a=\mp1$ when $\operatorname{Re}\tau=\pm\frac12$.

(a) For $\tau'=\gamma\tau\ne\tau$, the cases are $\tau'=\tau\pm1$ on the vertical sides, $\tau'=-1/\tau$ on the unit circle, and the cases with
$\tau\in\{\rho,-\overline\rho\}$, where $\tau'\in\{\rho,\rho+1\}$ is again a translate.

(b) $\gamma\tau=\tau$ with $c=0$ forces $\gamma=\pm I$. With $c=1$, $d=0$, $a=0$: $-1/\tau=\tau$ means $\tau=i$, and $\gamma=S$. The remaining fixers occur only at
$\tau=\rho$ or $-\overline\rho$. At $\rho$, using $\rho^2+\rho+1=0$: for $d=0$, $a=-1$, $\gamma\rho=-1-\overline\rho=\rho$; for $d=1$ (so $a-b=1$),
$\gamma\rho=\frac{a\rho+b}{\rho+1}=-\rho(a\rho+b)=a+\rho$, which is $\rho$ for $a=0$, $b=-1$, that is $\gamma=ST$. So the stabiliser of $\rho$ is
$\pm\{I,ST,(ST)^2\}$, where $(ST)^2=\bigl(\begin{smallmatrix}-1&-1\\1&0\end{smallmatrix}\bigr)$ and $(ST)^3=-I$. The point $\rho+1=T\rho$ has the conjugate
stabiliser $T\langle ST\rangle T^{-1}=\langle TS\rangle$.
\end{proof}

So $F$, with its vertical sides glued by $T$ and the two halves of its circular arc glued by $S$, is a picture of the orbit space $\mathrm{SL}_2(\Z)\backslash\mathbb H$:
a sphere with one point missing (the cusp at $i\infty$) and two special points, $i$ and $\rho$, where the gluing folds the surface onto itself.

\section{The valence formula}

For $f$ meromorphic and not identically zero near $P$, let $v_P(f)$ be its order at $P$; for $f\in M_k$, let $v_\infty(f)$ be the order at $q=0$ of its $q$-expansion.
If $f$ is modular, $v_{\gamma P}(f)=v_P(f)$, since $f(\gamma\tau)=(c\tau+d)^kf(\tau)$ and $c\tau+d\ne0$.

\begin{theorem}[valence formula]\label{thm:valence}
Let $f\in M_k$, $f\ne0$. Then
\[
v_\infty(f)+\frac12v_i(f)+\frac13v_\rho(f)+\sum_Pv_P(f)=\frac k{12},
\]
where the sum runs over the orbits of points $P$ of $\mathbb H$ other than those of $i$ and $\rho$.
\end{theorem}
\begin{proof}
Write $f(\tau)=h(q)$. As $h(q)q^{-v_\infty(f)}$ is non-zero at $q=0$, $f$ has no zeros with $\operatorname{Im}\tau>T_0$ for some $T_0$; and $F\cap\{\operatorname{Im}\tau\le T_0\}$ is
compact, so $f$ has finitely many zeros in $F$, and the sum is finite. Assume first that $f$ has no zeros on the boundary of $F$, except possibly at $i$, $\rho$ and
$\rho+1$. Let $T>T_0$, $\varepsilon>0$ small, and let $\mathcal C$ be the contour of Figure~\ref{fig:contour}: the boundary of $F\cap\{\operatorname{Im}\tau\le T\}$,
traversed counterclockwise, with small arcs of radius $\varepsilon$ around $\rho$, $i$ and $\rho+1$ replacing the corners, so that these points lie outside. By the
argument principle (the residue theorem applied to $f'/f$, Volume~4, Theorem~9.26),
\[
\frac1{2\pi i}\int_{\mathcal C}\frac{f'}f\,\mathrm d\tau=\sum_Pv_P(f),
\]
the sum over the orbits other than those of $i$ and $\rho$ (each such orbit meets the inside of $\mathcal C$ once, by Theorem~\ref{thm:funddomain}(a)). We compute the
left side piece by piece.

\emph{The top.} On $\operatorname{Im}\tau=T$, from $\frac12+iT$ to $-\frac12+iT$, $q=e^{2\pi i\tau}$ runs once clockwise around the circle $|q|=e^{-2\pi T}$, and
$\frac{f'}f\,\mathrm d\tau=\frac{h'}h\,\mathrm dq$. This piece gives $-v_\infty(f)$.

\emph{The vertical sides} are exchanged by $\tau\mapsto\tau+1$, which preserves $f$, and are traversed in opposite directions: they cancel.

\emph{The small arcs.} Near a zero $P$ of order $m$, $\frac{f'}f=\frac m{\tau-P}+(\text{holomorphic})$, so a clockwise arc of angle $\theta$ around $P$ contributes
$-\frac{m\theta}{2\pi}+O(\varepsilon)$. The interior angle of $F$ at $\rho$ and at $\rho+1$ is $\frac\pi3$ (between the vertical line and the unit circle, whose tangent at
$\rho$ makes an angle of $\frac\pi6$ with the horizontal), and at $i$ it is $\pi$. As $v_{\rho+1}(f)=v_\rho(f)$, these arcs give
$-\frac16v_\rho(f)-\frac16v_\rho(f)-\frac12v_i(f)+O(\varepsilon)$.

\emph{The unit circle.} Let $A$ be the arc of the unit circle from $\rho$ to $i$ and $A'$ the arc from $i$ to $\rho+1$ (shortened by the small arcs). The map
$S\tau=-1/\tau$ sends $e^{i\theta}$ to $e^{i(\pi-\theta)}$, so it maps $A'$ onto $A$ with the orientation reversed. Differentiating $f(S\tau)=\tau^kf(\tau)$ gives
\[
\frac{f'}f(S\tau)\,\mathrm d(S\tau)=k\frac{\mathrm d\tau}\tau+\frac{f'}f(\tau)\,\mathrm d\tau,
\]
so
\[
\int_A\frac{f'}f+\int_{A'}\frac{f'}f=\int_{A'}\Bigl(\frac{f'}f(\tau)\,\mathrm d\tau-\frac{f'}f(S\tau)\,\mathrm d(S\tau)\Bigr)=-k\int_{A'}\frac{\mathrm d\tau}\tau.
\]
On $A'$, $\tau=e^{i\theta}$ with $\theta$ going from about $\frac\pi2$ to about $\frac\pi3$, so $\int_{A'}\frac{\mathrm d\tau}\tau\to-\frac{i\pi}6$, and this piece gives
$\frac1{2\pi i}\cdot\frac{ik\pi}6=\frac k{12}$ in the limit.

Letting $\varepsilon\to0$ gives the formula. If $f$ has zeros on the boundary other than $i,\rho,\rho+1$, they come in pairs $P$, $P+1$ on the vertical sides and
$P$, $-1/P$ on the arc (Theorem~\ref{thm:funddomain}); deform the contour by small arcs that leave out one point of each pair and take in the other, symmetrically
under $T$ and $S$, and the same computation applies.
\end{proof}

\begin{figure}[ht]
\centering\fig{ch14-contour}
\caption{The contour for the valence formula: the boundary of the fundamental domain $F$, cut off at height $T$, with small arcs around $\rho$, $i$ and $-\overline\rho=\rho+1$.}
\label{fig:contour}
\end{figure}

\begin{think}
Use the valence formula to show again that $M_2=0$: which non-negative values can $v_\infty+\frac12v_i+\frac13v_\rho+\sum v_P$ take, and is $\frac16$ one of them?
\end{think}

The formula leaves very little room. For $k=4$, $\frac4{12}=\frac13$ forces $v_\rho=1$ and nothing else: \emph{$E_4$ vanishes exactly on the orbit of $\rho$}, simply. For
$k=6$, $\frac12$ forces $v_i=1$: \emph{$E_6$ vanishes exactly on the orbit of $i$}. For $k=12$, $\Delta$ has $v_\infty=1$ and no other zero, as we knew.

\begin{theorem}[the $j$-invariant]\label{thm:jinv}
The function
\[
j=\frac{E_4^3}\Delta=\frac{1728E_4^3}{E_4^3-E_6^2}=\frac1q+744+196884q+21493760q^2+\cdots
\]
is holomorphic on $\mathbb H$, invariant under $\mathrm{SL}_2(\Z)$, has a simple pole at $q=0$, and induces a bijection $\mathrm{SL}_2(\Z)\backslash\mathbb H\to\C$. Moreover
$j(\rho)=0$ and $j(i)=1728$.
\end{theorem}
\begin{proof}
$\Delta$ has no zeros, so $j$ is holomorphic, and of weight $12-12=0$. Its expansion is $q^{-1}(1+\cdots)$ because $\Delta=q+\cdots$ and $E_4^3=1+\cdots$. Let
$c\in\C$ and $f_c=E_4^3-c\Delta\in M_{12}$. Then $j(\tau)=c$ iff $f_c(\tau)=0$, and $v_\infty(f_c)=0$. The valence formula for $k=12$ reads
$\frac12v_i+\frac13v_\rho+\sum v_P=1$, whose only solutions in non-negative integers are $(v_i,v_\rho,\sum v_P)=(2,0,0)$, $(0,3,0)$ or $(0,0,1)$ with a single
$P$. In each case $f_c$ vanishes on exactly one orbit. So $j$ takes the value $c$ on exactly one orbit. Finally $E_4(\rho)=0$ gives $j(\rho)=0$, and $E_6(i)=0$ gives
$j(i)=\frac{1728E_4(i)^3}{E_4(i)^3}=1728$.
\end{proof}

Over $\C$, an elliptic curve is $\C/\Lambda$ for a lattice $\Lambda=\Z\omega_1+\Z\omega_2$ (Weierstrass), and two such curves are isomorphic iff the lattices are
homothetic, that is, iff $\tau=\omega_1/\omega_2$ and $\tau'$ are in the same orbit. So Theorem~\ref{thm:jinv} says that the number $j(\tau)$ classifies complex
elliptic curves up to isomorphism; for $y^2=x^3+ax+b$ it is $1728\frac{4a^3}{4a^3+27b^2}$. That $j$ has integer coefficients (Exercise~\ref{ex:jint}) has
startling consequences: for $\tau=\frac{1+\sqrt{-163}}2$, $j(\tau)=-640320^3$ is an integer (see the remark on complex multiplication in Section~\ref{rem:cm}), so $e^{\pi\sqrt{163}}=-q^{-1}$ is within $10^{-12}$ of the
integer $640320^3+744$.

\section{Hecke operators}

Volume~5 defined $T_p$ by a sum over matrices of determinant $p$. The cleanest way to define $T_n$ for all $n$ and to see how they multiply is to go back to lattices.

\subsection*{Modular forms as functions of lattices}

A \emph{lattice} is a subgroup $\Lambda=\Z\omega_1+\Z\omega_2$ of $\C$ with $\omega_1,\omega_2$ independent over $\R$; we order the basis so that
$\operatorname{Im}(\omega_1/\omega_2)>0$. A function $F$ on lattices has \emph{weight $k$} if $F(\lambda\Lambda)=\lambda^{-k}F(\Lambda)$ for $\lambda\in\C^\times$.

\begin{lemma}\label{lem:latticefn}
The maps $F\mapsto f$, $f(\tau)=F(\Z\tau+\Z)$, and $f\mapsto F$, $F(\Z\omega_1+\Z\omega_2)=\omega_2^{-k}f(\omega_1/\omega_2)$, are inverse bijections between
functions of weight $k$ on lattices and functions $f$ on $\mathbb H$ with $f|_k\gamma=f$ for all $\gamma\in\mathrm{SL}_2(\Z)$.
\end{lemma}
\begin{proof}
If $F$ has weight $k$ and $\gamma=\bigl(\begin{smallmatrix}a&b\\c&d\end{smallmatrix}\bigr)$, then $\Z(a\tau+b)+\Z(c\tau+d)=\Z\tau+\Z$, so
$f(\gamma\tau)=F\bigl((c\tau+d)^{-1}(\Z\tau+\Z)\bigr)=(c\tau+d)^kf(\tau)$. Conversely, two oriented bases of $\Lambda$ differ by some $\gamma\in\mathrm{SL}_2(\Z)$:
$(\omega_1',\omega_2')=(a\omega_1+b\omega_2,c\omega_1+d\omega_2)$. With $\tau=\omega_1/\omega_2$, $\omega_2'=\omega_2(c\tau+d)$ and
$\omega_2'^{-k}f(\gamma\tau)=\omega_2^{-k}(c\tau+d)^{-k}(c\tau+d)^kf(\tau)$, so $F$ is well defined; it has weight $k$, and the two maps are inverse.
\end{proof}

\begin{lemma}\label{lem:sublattices}
The sublattices of index $n$ of $\Lambda=\Z\omega_1+\Z\omega_2$ are the lattices $\Z(a\omega_1+b\omega_2)+\Z d\omega_2$ with $ad=n$, $a,d\ge1$ and $0\le b<d$, each
exactly once. In particular, $\Lambda$ has $\sigma_1(n)$ sublattices of index $n$.
\end{lemma}
\begin{proof}
Let $\Lambda'$ have index $n$. Then $\Lambda'\cap\Z\omega_2=\Z d\omega_2$ for some $d\ge1$, and the image of $\Lambda'$ in $\Lambda/\Z\omega_2\cong\Z\omega_1$ is $\Z a\omega_1$ for some
$a\ge1$; choose $a\omega_1+b\omega_2\in\Lambda'$. Then $\Lambda'=\Z(a\omega_1+b\omega_2)+\Z d\omega_2$, and its index is $ad$ (the determinant of the basis). Changing $b$
by a multiple of $d$ does not change $\Lambda'$, so we can take $0\le b<d$; and $a,d$ are determined by $\Lambda'$, and then $b\bmod d$ is too (the elements of
$\Lambda'$ with $\omega_1$-coordinate $a$ have $\omega_2$-coordinate $\equiv b\pmod d$).
\end{proof}

For $F$ of weight $k$ and $n\ge1$, let $T(n)F(\Lambda)=\sum_{[\Lambda:\Lambda']=n}F(\Lambda')$, a finite sum which is again of weight $k$, and $R(\lambda)F(\Lambda)=F(\lambda\Lambda)$. For
$f\in M_k$ define
\[
T_nf=n^{k-1}\,T(n)F,\qquad F\text{ the lattice function of }f.
\]
By Lemma~\ref{lem:sublattices}, with $\Lambda=\Z\tau+\Z$,
\begin{equation}\label{eq:Tn}
T_nf(\tau)=n^{k-1}\sum_{\substack{ad=n\\0\le b<d}}d^{-k}f\Bigl(\frac{a\tau+b}d\Bigr),
\end{equation}
since $F(\Z(a\tau+b)+\Z d)=d^{-k}F(\Z\frac{a\tau+b}d+\Z)$. For $n=p$ this is the operator $T_p$ of Volume~5.

\begin{theorem}[Hecke]\label{thm:heckeTn}
Let $k\ge4$ be even and $f=\sum_{m\ge0}c_mq^m\in M_k$.
\begin{enumerate}[label=(\alph*)]
\item $T_nf\in M_k$, with $q$-expansion coefficients
\[
c_m(T_nf)=\sum_{a\mid\gcd(m,n)}a^{k-1}c_{mn/a^2};
\]
in particular $c_0(T_nf)=\sigma_{k-1}(n)c_0$ and $c_1(T_nf)=c_n$. So $T_n$ maps $S_k$ to $S_k$.
\item $T_mT_n=T_{mn}$ if $\gcd(m,n)=1$, and $T_pT_{p^r}=T_{p^{r+1}}+p^{k-1}T_{p^{r-1}}$ for $p$ prime and $r\ge1$.
\item The operators $T_n$ commute with each other.
\end{enumerate}
\end{theorem}
\begin{proof}
(a) $T(n)F$ has weight $k$, so $T_nf|_k\gamma=T_nf$ by Lemma~\ref{lem:latticefn}, and $T_nf$ is holomorphic on $\mathbb H$ by \eqref{eq:Tn}. For fixed $a,d$,
$\sum_{b=0}^{d-1}e^{2\pi im(a\tau+b)/d}$ is $d\,q^{am/d}$ if $d\mid m$ and $0$ otherwise. So
\[
T_nf=n^{k-1}\sum_{ad=n}d^{1-k}\sum_{m'\ge0}c_{dm'}q^{am'}=\sum_{ad=n}a^{k-1}\sum_{m'\ge0}c_{dm'}q^{am'},
\]
holomorphic at infinity, and the coefficient of $q^m$ collects the terms with $am'=m$: $a\mid m$, $a\mid n$, $dm'=\frac na\cdot\frac ma=\frac{mn}{a^2}$.

(b) $T(m)T(n)F(\Lambda)$ is the sum of $F(\Lambda'')$ over the chains $\Lambda''\subseteq\Lambda'\subseteq\Lambda$ with $[\Lambda:\Lambda']=n$ and $[\Lambda':\Lambda'']=m$. If
$\gcd(m,n)=1$, each $\Lambda''$ of index $mn$ lies in exactly one chain: $\Lambda/\Lambda''$ is an abelian group of order $mn$, and $\Lambda'/\Lambda''$ must be its
unique subgroup of order $m$ (the elements killed by $m$). So $T(m)T(n)=T(mn)$, and $T_mT_n=T_{mn}$.

For $T(p)T(p^r)$, the chains are $\Lambda''\subseteq\Lambda'\subseteq\Lambda$ with $[\Lambda:\Lambda']=p$, $[\Lambda':\Lambda'']=p^r$. Fix $\Lambda''$ of index $p^{r+1}$ and
count the $\Lambda'$. A sublattice of index $p$ contains $p\Lambda$, and these correspond to the $p+1$ lines of $\Lambda/p\Lambda\cong(\Z/p)^2$; $\Lambda'$ contains
$\Lambda''$ iff its line contains the image of $\Lambda''$. If $\Lambda''\subseteq p\Lambda$, the image is $0$ and there are $p+1$ choices. Otherwise the image is
non-zero, and it is not all of $\Lambda/p\Lambda$: if $\Lambda''+p\Lambda=\Lambda$, the finite $p$-group $G=\Lambda/\Lambda''$ would satisfy $G=pG$, hence $G=0$ (as
$|pG|<|G|$ for a non-trivial $p$-group), but $|G|=p^{r+1}$. So the image is a line, and there is exactly one choice. Hence
\[
T(p)T(p^r)F(\Lambda)=T(p^{r+1})F(\Lambda)+p\sum_{\Lambda''\subseteq p\Lambda,\ [\Lambda:\Lambda'']=p^{r+1}}F(\Lambda'').
\]
The $\Lambda''$ in the last sum are the $p\Lambda'''$ with $[\Lambda:\Lambda''']=p^{r-1}$ (as $[\Lambda:p\Lambda]=p^2$), and $F(p\Lambda''')=p^{-k}F(\Lambda''')$. So
$T(p)T(p^r)=T(p^{r+1})+p^{1-k}T(p^{r-1})$, and multiplying by $p^{(r+1)(k-1)}$ gives the formula for $T_p$.

(c) By (b), each $T_{p^r}$ is a polynomial in $T_p$, and $T_n=\prod T_{p^{r_p}}$ over the prime powers dividing $n$. Operators $T_p,T_\ell$ for distinct primes commute,
since $T_pT_\ell=T_{p\ell}=T_\ell T_p$. So all $T_n$ lie in the commutative algebra generated by the $T_p$.
\end{proof}

A modular form $f\ne0$ is a \emph{Hecke eigenform} if $T_nf=\lambda_nf$ for all $n\ge1$.

\begin{corollary}\label{cor:eigen}
\begin{enumerate}[label=(\alph*)]
\item If $k\ge4$ and $S_k\ne0$, then $S_k$ contains a Hecke eigenform.
\item If $f=\sum c_mq^m\in S_k$ is a Hecke eigenform, then $c_1\ne0$; normalised by $c_1=1$, its coefficients are its eigenvalues, $c_n=\lambda_n$, and
\[
c_mc_n=\sum_{a\mid\gcd(m,n)}a^{k-1}c_{mn/a^2}.
\]
In particular $c_{mn}=c_mc_n$ for $\gcd(m,n)=1$ and $c_{p^{r+1}}=c_pc_{p^r}-p^{k-1}c_{p^{r-1}}$.
\end{enumerate}
\end{corollary}
\begin{proof}
(a) The operators $T_n$ on the finite-dimensional space $S_k$ span a finite-dimensional commutative algebra $\mathcal T$; pick generators $A_1,\dots,A_s$. An eigenspace
of $A_1$ is stable under the others (if $A_1v=\lambda v$ then $A_1A_iv=A_iA_1v=\lambda A_iv$); inside it, take an eigenspace of $A_2$, and so on. The last space is
non-zero and consists of common eigenvectors.

(b) By Theorem~\ref{thm:heckeTn}(a), $c_n=c_1(T_nf)=\lambda_nc_1$. If $c_1=0$, all $c_n$ vanish and $f=0$. With $c_1=1$, $\lambda_mc_n=c_n(T_mf)=\sum_{a\mid(m,n)}a^{k-1}c_{mn/a^2}$.
\end{proof}

When $\dim S_k=1$, that is for $k=12,16,18,20,22,26$, the only cusp form is an eigenform: $\Delta$, $\Delta E_4$, $\Delta E_6$, $\Delta E_8$, $\Delta E_{10}$, $\Delta E_{14}$ all have
multiplicative coefficients (Exercise~\ref{ex:deltae4}). For $k=24$, $S_{24}$ has dimension $2$, and $T_2$ has the characteristic polynomial $x^2-1080x-20468736$, whose
discriminant is $2^{6}3^2\cdot144169$: the two eigenforms have coefficients in $\Q(\sqrt{144169})$, and neither has rational coefficients (a direct computation gives them).

\begin{remark}[Petersson]
The \emph{Petersson inner product} on $S_k$ is
\[
\langle f,g\rangle=\int_Ff(\tau)\overline{g(\tau)}\,y^k\,\frac{\mathrm dx\,\mathrm dy}{y^2},
\]
and it makes every $T_n$ self-adjoint (Petersson, 1939). Hence $S_k$ has a basis of eigenforms, orthogonal, with real eigenvalues, and two normalised eigenforms with the same eigenvalues are equal (\emph{multiplicity one}).
We shall not need this; see Diamond--Shurman, Chapter~5, or Serre's \emph{Course in Arithmetic}, Chapter~VII.
\end{remark}

\section{$L$-functions and Hecke's functional equation}

\begin{proposition}[Hecke's bound]\label{prop:heckebound}
If $f=\sum c_nq^n\in S_k$, then $c_n=O(n^{k/2})$.
\end{proposition}
\begin{proof}
The function $\phi(\tau)=y^{k/2}|f(\tau)|$, $y=\operatorname{Im}\tau$, is invariant: $\operatorname{Im}(\gamma\tau)^{k/2}=y^{k/2}|c\tau+d|^{-k}$ and
$|f(\gamma\tau)|=|c\tau+d|^k|f(\tau)|$. On $F$, $f=q\,g(q)$ with $g$ bounded for $|q|\le e^{-\pi\sqrt3}$, so $\phi\le Cy^{k/2}e^{-2\pi y}$ is bounded on $F$, hence on $\mathbb H$,
say by $M$. Then for every $y>0$,
\[
|c_n|=\Bigl|\int_0^1f(x+iy)e^{-2\pi in(x+iy)}\,\mathrm dx\Bigr|\le My^{-k/2}e^{2\pi ny},
\]
and $y=\frac1n$ gives $|c_n|\le Me^{2\pi}n^{k/2}$.
\end{proof}

So $L(f,s)=\sum_{n\ge1}c_nn^{-s}$ converges absolutely for $\operatorname{Re}s>\frac k2+1$. For a normalised eigenform, Corollary~\ref{cor:eigen}(b) gives an Euler product
there, exactly as for $\tau$ in Volume~5:
\[
L(f,s)=\prod_p\frac1{1-c_pp^{-s}+p^{k-1-2s}}.
\]
(Expand each factor as $\sum_rc_{p^r}p^{-rs}$: the recursion $c_{p^{r+1}}=c_pc_{p^r}-p^{k-1}c_{p^{r-1}}$ says that $(1-c_pX+p^{k-1}X^2)\sum c_{p^r}X^r=1$.)

Hecke's theorem needs only one symmetry, and it is convenient to prove it at once in the generality needed for the last section.

\begin{theorem}[Hecke, 1936]\label{thm:heckefe}
Let $N\ge1$, $k\ge2$ even, and $f(\tau)=\sum_{n\ge1}c_nq^n$ holomorphic on $\mathbb H$, with $c_n=O(n^A)$ for some $A$. Assume that for some $\varepsilon=\pm1$,
\begin{equation}\label{eq:fricke}
f\Bigl(-\frac1{N\tau}\Bigr)=\varepsilon N^{k/2}\tau^kf(\tau)\qquad(\tau\in\mathbb H).
\end{equation}
Then
\[
\Lambda(s)=\Bigl(\frac{\sqrt N}{2\pi}\Bigr)^s\Gamma(s)L(f,s),\qquad L(f,s)=\sum_{n\ge1}\frac{c_n}{n^s}\quad(\operatorname{Re}s>A+1),
\]
extends to an entire function, bounded in every vertical strip of finite width, and
\[
\Lambda(k-s)=\varepsilon(-1)^{k/2}\Lambda(s).
\]
\end{theorem}
\begin{proof}
\emph{The Mellin transform.} For $n\ge1$ and $\operatorname{Re}s>0$, substituting $t=2\pi ny$ gives $\int_0^\infty e^{-2\pi ny}y^{s-1}\,\mathrm dy=(2\pi n)^{-s}\Gamma(s)$. For
$\sigma=\operatorname{Re}s>A+1$, $\sum_n|c_n|\int_0^\infty e^{-2\pi ny}y^{\sigma-1}\,\mathrm dy=\Gamma(\sigma)(2\pi)^{-\sigma}\sum|c_n|n^{-\sigma}<\infty$, so we may sum under the
integral:
\[
\Lambda(s)=N^{s/2}\int_0^\infty f(iy)\,y^{s-1}\,\mathrm dy.
\]
\emph{Decay.} For $y\ge y_0>0$,
\[
|f(iy)|\le e^{-2\pi y}\sum_n|c_n|e^{-2\pi(n-1)y_0}=C_{y_0}e^{-2\pi y}.
\]
So for $y_0=1/\sqrt N$, the integral
$I(s)=\int_{y_0}^\infty f(iy)y^{s-1}\,\mathrm dy$ converges for every $s$, uniformly on vertical strips, and defines an entire function, bounded on each strip.

\emph{The symmetry.} With $\tau=iy$, \eqref{eq:fricke} reads $f\bigl(\frac i{Ny}\bigr)=\varepsilon i^kN^{k/2}y^kf(iy)$. In $\int_0^{y_0}f(iy)y^{s-1}\,\mathrm dy$ substitute $y=\frac1{Nu}$,
so that $u$ runs from $\infty$ down to $y_0$ and $\mathrm dy=-\frac{\mathrm du}{Nu^2}$:
\[
\int_0^{y_0}f(iy)y^{s-1}\,\mathrm dy=\int_{y_0}^\infty\varepsilon i^kN^{k/2}u^kf(iu)\,(Nu)^{1-s}\frac{\mathrm du}{Nu^2}=\varepsilon i^kN^{k/2-s}I(k-s).
\]
Hence, for $\operatorname{Re}s>A+1$,
\[
\Lambda(s)=N^{s/2}I(s)+\varepsilon i^kN^{(k-s)/2}I(k-s).
\]
The right side is entire and bounded on vertical strips, which gives the continuation. Replacing $s$ by $k-s$ and using $(\varepsilon i^k)^2=1$ ($k$ is even) gives
$\Lambda(k-s)=\varepsilon i^k\Lambda(s)$, and $i^k=(-1)^{k/2}$.
\end{proof}

\begin{corollary}\label{cor:level1fe}
For $f\in S_k$, $k\ge4$, the function $\Lambda(s)=(2\pi)^{-s}\Gamma(s)L(f,s)$ is entire and $\Lambda(k-s)=(-1)^{k/2}\Lambda(s)$. In particular $L(\Delta,s)$ extends to an entire
function, with $(2\pi)^{-s}\Gamma(s)L(\Delta,s)$ invariant under $s\mapsto12-s$.
\end{corollary}
\begin{proof}
Take $N=1$: $f(-1/\tau)=\tau^kf(\tau)$, so $\varepsilon=1$, and $c_n=O(n^{k/2})$ by Proposition~\ref{prop:heckebound}.
\end{proof}

Compare with the Riemann zeta function (Volume~4, Chapter~10): there, $\Gamma(\frac s2)\pi^{-s/2}\zeta(s)$ is symmetric under $s\mapsto1-s$ and has poles at $0$ and $1$, which come from the constant
term of the theta function. A cusp form has no constant term, and its $L$-function is entire. Its critical line is $\operatorname{Re}s=\frac k2$, and the analogue of
the Riemann hypothesis, that all the non-trivial zeros lie on it, is open for every single cusp form; Figure~\ref{fig:hardy} shows the first zeros of $L(\Delta,s)$ on
$\operatorname{Re}s=6$.

\begin{figure}[ht]
\centering\fig{ch14-hardy}
\caption{The real functions $Z(t)=\pm|L(\frac k2+it)|$ (with the sign of $\Lambda(\frac k2+it)$, which is real) for $L(\Delta,s)$, $k=12$ (top), and for $L(f,s)$,
$f=\eta(\tau)^2\eta(11\tau)^2$, $k=2$ (bottom), $0\le t\le30$. Each sign change is a zero on the critical line. The first zero of $L(\Delta,s)$ is at $t=9.2224$;
the first of $L(f,s)$ at $t=6.3626$.}
\label{fig:hardy}
\end{figure}

\section{An elliptic curve and a modular form}

The functions with the symmetry \eqref{eq:fricke} for $N>1$ are modular forms for the group
\[
\Gamma_0(N)=\Bigl\{\begin{pmatrix}a&b\\c&d\end{pmatrix}\in\mathrm{SL}_2(\Z):\ N\mid c\Bigr\},
\]
which are \emph{eigenvectors of the Fricke involution} $\tau\mapsto-1/(N\tau)$; modular forms of \emph{level $N$}. We do not develop the theory of these forms, but
we can handle one of them completely with what we have.

\begin{proposition}\label{prop:eta11}
Let $f(\tau)=\eta(\tau)^2\eta(11\tau)^2=q\prod_{n\ge1}(1-q^n)^2(1-q^{11n})^2=\sum_{n\ge1}c_nq^n$. Then $c_n=O(n^2)$ and
\[
f\Bigl(-\frac1{11\tau}\Bigr)=-11\,\tau^2f(\tau).
\]
Consequently $\Lambda(s)=\bigl(\frac{\sqrt{11}}{2\pi}\bigr)^s\Gamma(s)L(f,s)$ extends to an entire function with $\Lambda(2-s)=\Lambda(s)$.
\end{proposition}
\begin{proof}
The exponent of $q$ in $\eta(\tau)^2\eta(11\tau)^2$ is $\frac{2+22}{24}=1$, which gives the product. By Volume~5, Theorem~12.7, applied at $11\tau$ and at $\tau$,
\[
\eta\Bigl(-\frac1{11\tau}\Bigr)^2=-11i\tau\,\eta(11\tau)^2,\qquad\eta\Bigl(-\frac1\tau\Bigr)^2=-i\tau\,\eta(\tau)^2,
\]
and the product is $(-i)^2\cdot11\tau^2f(\tau)=-11\tau^2f(\tau)$. For the growth, by Euler's pentagonal theorem $\prod(1-q^n)=\sum_j\pm q^{P_j}$ with distinct pentagonal
numbers $P_j=\frac{j(3j-1)}2$, at most $2\sqrt n+1$ of them $\le n$. So $c_n$, the coefficient of $q^{n-1}$ in the product of two copies of this series and two copies of
the same series in $q^{11}$, is a sum of at most $(2\sqrt n+1)^4$ terms $\pm1$. Theorem~\ref{thm:heckefe} applies with $N=11$, $k=2$, $\varepsilon=-1$, and
$\varepsilon(-1)^{k/2}=1$.
\end{proof}

Now compare the coefficients $c_p$ with the elliptic curve $E:y^2+y=x^3-x^2$, which has discriminant $-11$ (complete the square: $(2y+1)^2=4x^3-4x^2+1$) and good
reduction at every $p\ne11$. Let $N_p$ be its number of projective points over $\mathbb F_p$, and $a_p=p+1-N_p$.

\begin{center}
\begin{tabular}{c|ccccc}
$p$&$2$&$3$&$5$&$7$&$13$\\\hline
$N_p$&$5$&$5$&$5$&$10$&$10$\\
$a_p$&$-2$&$-1$&$1$&$-2$&$4$\\
$c_p$&$-2$&$-1$&$1$&$-2$&$4$
\end{tabular}\par\medskip
\begin{tabular}{c|ccccc}
$p$&$17$&$19$&$23$&$29$&$31$\\\hline
$N_p$&$20$&$20$&$25$&$30$&$25$\\
$a_p$&$-2$&$0$&$-1$&$0$&$7$\\
$c_p$&$-2$&$0$&$-1$&$0$&$7$
\end{tabular}
\end{center}

A direct computation verifies $a_p=c_p$ for all $p<10^5$, $p\ne11$. This is not a coincidence.

\begin{remark}[Eichler, Shimura, and the modularity theorem]\label{rem:modularity}
Eichler (1954) and Shimura (1958) showed that for a weight-$2$ eigenform $f$ on $\Gamma_0(N)$ with rational coefficients, the numbers $c_p$ are $p+1-N_p$ for an
elliptic curve over $\Q$ built from the modular curve $X_0(N)=\Gamma_0(N)\backslash\mathbb H\cup\{\text{cusps}\}$; for $N=11$ one gets a curve of conductor $11$, and every
such curve has the same $a_p$ as $E$. The \emph{modularity theorem} (Wiles, Taylor--Wiles 1995; Breuil, Conrad, Diamond, Taylor 2001) is the converse: for every elliptic
curve $E$ over $\Q$ of conductor $N$, there is such an $f$ with $c_p=a_p$ for all $p\nmid N$. Consequently $L(E,s)=\prod_{p\nmid N}(1-a_pp^{-s}+p^{1-2s})^{-1}\times(\text{bad factors})$
is the $L(f,s)$ of Theorem~\ref{thm:heckefe}: it continues to $\C$ and satisfies $\Lambda(2-s)=\pm\Lambda(s)$. Before 1995 this was not known for most curves, and without it
the Birch--Swinnerton-Dyer conjecture, which is about $L(E,1)$, could not even be stated. For our $E$, $L(E,1)=0.2538\ldots\ne0$, and indeed $E(\Q)$ is finite: it is
$\Z/5$, generated by $(0,0)$.
\end{remark}

\begin{remark}[Fermat]
If $a^\ell+b^\ell=c^\ell$ with $\ell\ge5$ prime and $abc\ne0$, Frey's curve $y^2=x(x-a^\ell)(x+b^\ell)$ is semistable after normalisation; by Wiles it is modular, of some level $N$;
by Ribet's theorem (1986, conjectured by Serre), its mod-$\ell$ representation comes from a weight-$2$ cusp form of level $2$. There is none: $\Gamma_0(2)$ has no non-zero cusp
form of weight $2$, by a dimension count analogous to Volume~5, Theorem~12.9. So Fermat's Last Theorem ends, like Ramanujan's congruence, with a space of modular forms
that is too small.
\end{remark}

\begin{remark}[complex multiplication]\label{rem:cm}
When $\tau$ is a root of a quadratic equation, the curve $\C/(\Z\tau+\Z)$ has extra endomorphisms, and $j(\tau)$ is an algebraic integer which generates the Hilbert
class field of $\Q(\tau)$ (Chapter~13): this is Kronecker's \emph{Jugendtraum}, the class field theory of imaginary quadratic fields. When $\Q(\tau)$ has class number $1$,
$j(\tau)\in\Z$. The largest such field is $\Q(\sqrt{-163})$, and $j\bigl(\frac{1+\sqrt{-163}}2\bigr)=-640320^3$. The Chudnovsky brothers' series for $\frac1\pi$, used for
most record computations of $\pi$ since 1989, comes from this value (a direct computation gives $10^5$ digits of $\pi$ with it).
\end{remark}

\begin{beauty}
Two numbers, $c_p$ and $a_p$, one computed from a product of infinitely many factors $(1-q^n)$, the other by counting solutions of a cubic congruence, agree for every
prime that has been checked. Behind the coincidence is the whole of this book's last third: the symmetry $\tau\mapsto-1/(11\tau)$ gives an $L$-function with a functional equation, and that equation is proved. Hasse's bound, proved in Chapter~16 for curves $y^2=f(x)$, gives $|a_p|\le2\sqrt p$ for this curve. The equality $c_p=a_p$, and with it $|c_p|\le2\sqrt p$, is a numerical check for $p<10^5$ together with the Eichler--Shimura theorem, quoted in the remark above and not proved here. The arithmetic of the curve, its rank and its torsion, is conjecturally read off at the single point
$s=1$. Fermat's Last Theorem was one consequence of making the coincidence a theorem.
\end{beauty}

\section{Exercises}

\begin{exercise}[$\star$]\label{ex:jvalues}
Using the valence formula, show that $E_4$ has no zeros in $F$ other than $\rho$ and $\rho+1$, and $E_6$ none other than $i$. Deduce $j(\rho)=0$ and $j(i)=1728$, and show
that $j$ is real on the boundary of $F$ and on the imaginary axis.
\end{exercise}

\begin{exercise}[$\star$]\label{ex:eisenhecke}
Show, using the coefficient formula of Theorem~\ref{thm:heckeTn}, that $T_nE_k=\sigma_{k-1}(n)E_k$. Deduce that the Dirichlet series of the coefficients of
$E_k$, normalised by $c_1=1$, is $\sum_{n\ge1}\sigma_{k-1}(n)n^{-s}=\zeta(s)\zeta(s-k+1)$, with an Euler product of the same shape as for cusp forms.
\end{exercise}

\begin{exercise}[$\star$]\label{ex:trivialzeros}
Show that $L(\Delta,s)$ vanishes at $s=0,-1,-2,\dots$, and that its other zeros lie in the strip $5\le\operatorname{Re}s\le7$. (Use the
Euler product, Hecke's bound and the functional equation.)
\end{exercise}

\begin{exercise}[$\star\star$]\label{ex:deltae4}
Show that $\Delta E_4=q+216q^2-3348q^3+13888q^4+\cdots$ is a normalised eigenform, and check $c_6=c_2c_3$ and $c_4=c_2^2-2^{15}$ on the coefficients.
\end{exercise}

\begin{exercise}[$\star\star$]\label{ex:jint}
Show that $j$ has integer coefficients, and compute $j=q^{-1}+744+196884q+\cdots$. Show that $196884=196883+1$, where $196883$ is the smallest dimension of a non-trivial
representation of the Monster group (this last fact is to be admitted; it is the starting point of ``monstrous moonshine'').
\end{exercise}

\begin{exercise}[$\star\star$]\label{ex:level2}
Let $g(\tau)=\eta(\tau)^8\eta(2\tau)^8=q-8q^2+12q^3+64q^4-\cdots$. Show that $g(-\frac1{2\tau})=16\tau^8g(\tau)$, and deduce the analytic continuation and functional equation
of $L(g,s)$. What is the sign?
\end{exercise}

\begin{exercise}[$\star\star\star$]\label{ex:e8theta}
Let $\Lambda_8\subset\R^8$ be the $E_8$ lattice: the vectors $x$ with all $x_i\in\Z$ or all $x_i\in\Z+\frac12$, and $\sum x_i\in2\Z$. Show that $x\cdot x\in2\Z$ for
$x\in\Lambda_8$, that $\Lambda_8$ has covolume $1$ and equals its dual. Using Poisson summation (Volume~4, Theorem~8.9, in eight variables), show that
$\theta(\tau)=\sum_{x\in\Lambda_8}q^{x\cdot x/2}$ satisfies $\theta(-1/\tau)=\tau^4\theta(\tau)$. Deduce $\theta=E_4$, so that $\Lambda_8$ has $240\sigma_3(n)$ vectors with
$x\cdot x=2n$; in particular $240$ vectors of length $\sqrt2$.
\end{exercise}

\section{Hints}

\begin{hint}{ex:jvalues}
$\frac4{12}=\frac13$, $\frac6{12}=\frac12$. For reality: $\overline{f(\tau)}=f(-\overline\tau)$ when $f$ has real coefficients.
\end{hint}
\begin{hint}{ex:eisenhecke}
The coefficient $c_m(T_nE_k)$ involves $\sum_{a\mid(m,n)}a^{k-1}\sigma_{k-1}(mn/a^2)$. Prove the identity
\[
\sigma_{k-1}(m)\,\sigma_{k-1}(n)=\sum_{a\mid(m,n)}a^{k-1}\sigma_{k-1}\Bigl(\frac{mn}{a^2}\Bigr)
\]
for prime powers first.
\end{hint}
\begin{hint}{ex:trivialzeros}
$\Gamma$ has poles at $0,-1,-2,\dots$, and $\Lambda$ is entire. A convergent Euler product does not vanish.
\end{hint}
\begin{hint}{ex:deltae4}
$\dim S_{16}=1$. $c_4=c_2^2-2^{15}$ is the recursion with $k-1=15$.
\end{hint}
\begin{hint}{ex:jint}
$E_4^3\in1+240q\Z[[q]]$ and $\Delta^{-1}=q^{-1}\prod(1-q^n)^{-24}\in q^{-1}\Z[[q]]$.
\end{hint}
\begin{hint}{ex:level2}
Apply Dedekind's formula at $2\tau$ and at $\tau$; then Theorem~\ref{thm:heckefe} with $N=2$, $k=8$, and Euler's pentagonal theorem for the growth.
\end{hint}
\begin{hint}{ex:e8theta}
The Fourier transform of $e^{\pi i\tau x\cdot x}$ on $\R^8$ is $(-i\tau)^{-4}e^{\pi i(-1/\tau)y\cdot y}$. For self-duality, show $\Lambda_8\subseteq\Lambda_8^*$ and compare covolumes.
\end{hint}

\section{Solutions}

\begin{solution}{ex:jvalues}
\idea The valence formula with a small right-hand side has only one solution.
\writeup For $f=E_4$, $v_\infty=0$ (constant term $1$), and $\frac12v_i+\frac13v_\rho+\sum v_P=\frac13$ forces $v_\rho=1$ and everything else $0$ (any other term is at least
$\frac12$, and $v_\rho\ge2$ gives $\frac23$). Similarly for $E_6$: $\frac12$ forces $v_i=1$. Then $j(\rho)=\frac{E_4(\rho)^3}{\Delta(\rho)}=0$, and since
$1728\Delta=E_4^3-E_6^2$, $j(i)=\frac{1728E_4(i)^3}{E_4(i)^3}=1728$, with $E_4(i)\ne0$.

$j$ has real coefficients, so $\overline{j(\tau)}=j(-\overline\tau)$ (as $\overline{e^{2\pi i\tau}}=e^{2\pi i(-\overline\tau)}$). On the imaginary axis and on
$\operatorname{Re}\tau=-\frac12$ (where $-\overline\tau=\tau+1$), $j(-\overline\tau)=j(\tau)$; on the unit circle, $-\overline\tau=-1/\tau$ and again $j(-\overline\tau)=j(\tau)$. So $j$ is real
there. In fact $j$ maps the left half of $F$ bijectively onto the upper or lower half-plane, and its boundary onto $\R$: $(-\infty,0]$ on the left side, $[0,1728]$ on the
arc, $[1728,\infty)$ on the imaginary axis.
\end{solution}

\begin{solution}{ex:eisenhecke}
\idea A multiplicativity identity for $\sigma_{k-1}$, checked on prime powers.
\writeup Write $s=\sigma_{k-1}$. By Theorem~\ref{thm:heckeTn}(a), with $E_k=1+\kappa\sum s(m)q^m$, $\kappa=-\frac{2k}{B_k}$: $c_0(T_nE_k)=s(n)$ and, for $m\ge1$,
$c_m(T_nE_k)=\kappa\sum_{a\mid(m,n)}a^{k-1}s(mn/a^2)$. So it suffices to prove $s(m)s(n)=\sum_{a\mid(m,n)}a^{k-1}s(mn/a^2)$. Both sides are multiplicative in the pair
$(m,n)$ (if $m=m_1m_2$, $n=n_1n_2$ with $m_1n_1$ prime to $m_2n_2$, the divisors $a$ split accordingly), so take $m=p^i$, $n=p^j$, $i\le j$, and $x=p^{k-1}$:
\[
\sum_{t=0}^ix^t\cdot s(p^{i+j-2t})=\sum_{t=0}^ix^t\sum_{u=0}^{i+j-2t}x^u ,
\]
which counts each $x^w$ ($0\le w\le i+j$) once for each $t$ with $0\le t\le\min(i,w,i+j-w)$; and $s(p^i)s(p^j)=\sum_{u\le i}\sum_{v\le j}x^{u+v}$ counts $x^w$ the same number of
times, $\min(i,w,i+j-w)+1$. So $T_nE_k=s(n)E_k$. Normalised by $c_1=1$, the coefficients of $E_k$ are $\sigma_{k-1}(n)$ for $n\ge1$, and
$\sum\sigma_{k-1}(n)n^{-s}=\zeta(s)\zeta(s-k+1)$, the Dirichlet convolution of $1$ and $n^{k-1}$; its Euler factor at $p$ is
$\frac1{(1-p^{-s})(1-p^{k-1-s})}=\frac1{1-\sigma_{k-1}(p)p^{-s}+p^{k-1-2s}}$.
\end{solution}

\begin{solution}{ex:trivialzeros}
\idea $\Lambda$ is entire, while $\Gamma$ has poles; the Euler product converges to the right of the strip, and the functional equation reflects it.
\writeup $L(\Delta,s)=(2\pi)^s\Lambda(s)/\Gamma(s)$, and $1/\Gamma$ is entire with zeros exactly at $0,-1,-2,\dots$ (Volume~4, Theorem~9.36). So $L(\Delta,s)$ is entire and vanishes
at these points unless $\Lambda$ has a pole there, which it does not. For $\sigma=\operatorname{Re}s>7$, $\sum|\tau(n)|n^{-\sigma}<\infty$ by Hecke's bound $\tau(n)=O(n^6)$, so the Euler
product converges absolutely; an absolutely convergent product of non-zero factors is non-zero, so $L(\Delta,s)\ne0$, and $\Lambda(s)\ne0$ ($\Gamma$ has no zeros). For
$\operatorname{Re}s<5$, $\Lambda(s)=\Lambda(12-s)\ne0$, so $L(\Delta,s)=0$ only where $1/\Gamma(s)=0$. (Deligne's bound $|\tau(n)|\le d(n)n^{11/2}$ narrows the strip to
$\frac{11}2\le\operatorname{Re}s\le\frac{13}2$, and the analogue of the Riemann hypothesis says $\operatorname{Re}s=6$.)
\end{solution}

\begin{solution}{ex:deltae4}
\idea A one-dimensional space stable under all $T_n$ consists of eigenvectors.
\writeup $\Delta E_4\in S_{16}$, which has dimension $\dim M_4=1$ (Volume~5, Theorem~12.9(b)). As $T_nS_{16}\subseteq S_{16}$, $T_n(\Delta E_4)=\lambda_n\Delta E_4$, and its $q$-coefficient is
$1$. Multiplying $q-24q^2+252q^3-1472q^4$ by $1+240q+2160q^2+6720q^3$ gives $q+216q^2-3348q^3+13888q^4$. Corollary~\ref{cor:eigen} gives $c_4=c_2^2-2^{15}=46656-32768=13888$, as
found. Two more terms of each factor give $c_5=52110$ and $c_6=-723168=216\cdot(-3348)=c_2c_3$.
\end{solution}

\begin{solution}{ex:jint}
\idea Integrality of each factor of $E_4^3\cdot\Delta^{-1}$.
\writeup $E_4=1+240\sum\sigma_3(n)q^n$ has integer coefficients, and so does $E_4^3$. The product $\prod(1-q^n)^{24}$ has integer coefficients and constant term $1$, so its
inverse in $\Z[[q]]$ exists (solve for the coefficients one at a time). So $j=q^{-1}E_4^3\prod(1-q^n)^{-24}\in q^{-1}\Z[[q]]$. With $E_4^3=1+720q+179280q^2+\cdots$ and
$\prod(1-q^n)^{-24}=1+24q+324q^2+\cdots$, $j=q^{-1}(1+744q+(179280+720\cdot24+324)q^2+\cdots)=q^{-1}+744+196884q+\cdots$.
\end{solution}

\begin{solution}{ex:level2}
\idea Same computation as Proposition~\ref{prop:eta11}.
\writeup $\eta(-\frac1{2\tau})^8=(-2i\tau)^4\eta(2\tau)^8=16\tau^4\eta(2\tau)^8$ and $\eta(-\frac1\tau)^8=(-i\tau)^4\eta(\tau)^8=\tau^4\eta(\tau)^8$, so $g(-\frac1{2\tau})=16\tau^8g(\tau)=2^{8/2}\tau^8g(\tau)$,
which is \eqref{eq:fricke} with $N=2$, $k=8$, $\varepsilon=1$. The exponent of $q$ is $\frac{8+16}{24}=1$, and the coefficients are $O(n^8)$ by the pentagonal theorem (a sum of at most
$(2\sqrt n+1)^{16}$ terms $\pm1$). So $\Lambda(s)=(\frac{\sqrt2}{2\pi})^s\Gamma(s)L(g,s)$ is entire with $\Lambda(8-s)=(-1)^4\Lambda(s)=\Lambda(s)$: the sign is $+1$.
\end{solution}

\begin{solution}{ex:e8theta}
\idea Even, unimodular, so its theta series is a modular form of weight $4$, and $M_4$ is a line.
\writeup \emph{Even.} If all $x_i\in\Z$, $x\cdot x\equiv\sum x_i\equiv0\pmod2$. If all $x_i=\frac12+n_i$, $x\cdot x=\sum(n_i^2+n_i)+2$, and $n_i^2+n_i$ is even. The lattice
contains $D_8=\{x\in\Z^8:\sum x_i\text{ even}\}$, of covolume $2$, with index $2$ (the second coset is $D_8+(\frac12,\dots,\frac12)$), so its covolume is $1$.
\emph{Self-dual.} For $x,y\in\Lambda_8$, $x\cdot y=\frac12((x+y)\cdot(x+y)-x\cdot x-y\cdot y)\in\Z$, so $\Lambda_8\subseteq\Lambda_8^*$; both have covolume $1$, so they are equal.
\emph{Poisson.} For a lattice $L=M\Z^8$ of covolume $|\det M|$ and $g$ a Gaussian, applying the one-variable formula in each coordinate to $g(Mx)$ gives
$\sum_{x\in L}g(x)=\frac1{\operatorname{covol}L}\sum_{y\in L^*}\hat g(y)$. With $g(x)=e^{\pi i\tau x\cdot x}$, for $\tau=iy$ first and then by analytic continuation,
$\hat g(\xi)=(-i\tau)^{-4}e^{\pi i(-1/\tau)\xi\cdot\xi}$. So $\theta(\tau)=(-i\tau)^{-4}\theta(-1/\tau)$, that is $\theta(-1/\tau)=\tau^4\theta(\tau)$. Also $\theta(\tau+1)=\theta(\tau)$
because $x\cdot x/2\in\Z$, and $\theta$ is holomorphic at infinity with constant term $1$. So $\theta\in M_4=\C E_4$, and $\theta=E_4$.
\end{solution}

\section*{Sources and further reading}

\begin{itemize}
\item J.-P.~Serre, \emph{A Course in Arithmetic}, Springer GTM~7, 1973, Chapter~VII. The fundamental domain, the valence formula and the Hecke operators via lattices,
followed here.
\item F.~Diamond and J.~Shurman, \emph{A First Course in Modular Forms}, Springer GTM~228, 2005. Congruence subgroups, Petersson's inner product, Eichler--Shimura, and the
modularity theorem.
\item T.~M.~Apostol, \emph{Modular Functions and Dirichlet Series in Number Theory}, 2nd ed., Springer GTM~41, 1990. The $j$-function, its integrality, and Hecke's theory.
\item E.~Hecke, ``\"Uber die Bestimmung Dirichletscher Reihen durch ihre Funktionalgleichung'', \emph{Math.\ Ann.} 112 (1936), 664--699.
\item M.~Eichler, ``Quatern\"are quadratische Formen und die Riemannsche Vermutung f\"ur die Kongruenzzetafunktion'', \emph{Arch.\ Math.} 5 (1954), 355--366.
\item A.~Weil, ``\"Uber die Bestimmung Dirichletscher Reihen durch Funktionalgleichungen'', \emph{Math.\ Ann.} 168 (1967), 149--156.
\item Y.~Taniyama's problems from the 1955 Tokyo--Nikk\=o symposium, and the history of the conjecture: S.~Lang, ``Some history of the Shimura--Taniyama conjecture'',
\emph{Notices AMS} 42 (1995), 1301--1307.
\item G.~Frey, ``Links between stable elliptic curves and certain Diophantine equations'', \emph{Ann.\ Univ.\ Sarav.} 1 (1986); K.~Ribet, ``On modular representations of
$\mathrm{Gal}(\overline{\Q}/\Q)$ arising from modular forms'', \emph{Invent.\ Math.} 100 (1990), 431--476.
\item A.~Wiles, ``Modular elliptic curves and Fermat's Last Theorem'', \emph{Ann.\ of Math.} 141 (1995), 443--551; R.~Taylor and A.~Wiles, ``Ring-theoretic properties of certain
Hecke algebras'', same volume, 553--572; C.~Breuil, B.~Conrad, F.~Diamond and R.~Taylor, ``On the modularity of elliptic curves over $\Q$'', \emph{J.\ AMS} 14 (2001), 843--939.
\item D.~V.~Chudnovsky and G.~V.~Chudnovsky, ``Approximations and complex multiplication according to Ramanujan'', in \emph{Ramanujan Revisited}, Academic Press, 1988.
\end{itemize}
